\section{Three Changes, and One Status Code}
\label{sec:ch16-three-changes-and-one-status-code}

\index{breaking change!what a client actually receives|(}%
Section~\ref{sec:ch16-the-change-this-book-already-shipped} is a change that
answers correctly, which makes it the hardest of the three to feel. So take
three edits to the Sessions document, one at a time, compose each into a
config and point a router at it. The four services behind that router are
untouched in every run, and that is worth stopping on: for two of the three
edits, the router refuses the request against its own copy of the composed
schema and never asks a service anything. The whole failure lives in a
document.

\subsection{The Deprecated Field, Removed}

\index{deprecation!the removal chapter 4 refused to make}%
Chapter~\ref{ch:schema-design} deprecated \code{sessionById} and said that
deleting it instead would break every client still asking, on the deployment,
with no warning. Twelve chapters later the field is still there and the reason
to remove it has only got better. Take it out of the document, with its
\code{@deprecated} block, and compose:

\begin{minted}[fontsize=\footnotesize,breakanywhere]{text}
Router execution config successfully written to ".../router.json".
\end{minted}

\index{composition!approves a field removal}%
Exit zero. Nothing else printed, and nothing else to print: the field belonged
to one subgraph, no \code{@key} named it, no \code{@requires} named it, and
the graph merges exactly as it did before.

\index{deprecation!the client that did not read the reason}%
Now a client that never read the reason string, sending what it has sent since
chapter~\ref{ch:schema-design}:

\begin{minted}{graphql}
{ sessionById(id: 1) { title } }
\end{minted}

\begin{minted}[fontsize=\footnotesize,breakanywhere]{json}
{"errors":[{"message":"Cannot query field \"sessionById\" on type \"Query\".","path":["query"]}]}
\end{minted}

\index{HTTP 200!a removed field is still a 200}%
Two hundred. Not a four hundred, not a four oh four. An \code{errors} array,
no \code{data} key at all, and a status line that a client library will hand
back as success.

\subsection{The Argument Nobody Had Yesterday}

\index{breaking change!a required argument added}%
The second edit adds a required argument. \code{sessions} takes four paging
arguments and none of them is mandatory; give it a fifth, \code{track:
String!}, with no default, as a team would while adding a filter they consider
obviously necessary. It composes at exit zero, and the composed schema carries
\code{sessions(track: String!, first: Int, after: String, last: Int, before:
String)}.

\index{breaking change!the query this book has printed since chapter 3}%
The query that breaks is the one this book has printed since
chapter~\ref{ch:data-without-n-plus-1}:

\begin{minted}{graphql}
{ sessions(first: 10) { nodes { title } } }
\end{minted}

\begin{minted}[fontsize=\footnotesize,breakanywhere]{json}
{"errors":[{"message":"argument: track is required on field: sessions but missing","path":["query"]}]}
\end{minted}

\index{HTTP 200!a required argument is still a 200}%
Two hundred again, and it is worth putting this beside the misspelled field
chapter~\ref{ch:hot-chocolate-zero} sent. That one came back from Hot
Chocolate with a \code{locations} array, an \code{extensions} object and a
link to the clause of the specification it was enforcing. This one has a
\code{path} and nothing else.

\index{HTTP 200!the same failure was a 400 before the router}%
And the status code moved. Chapter~\ref{ch:hot-chocolate-zero}'s validation
error answered four hundred, which is that chapter's own measurement. Send the
same class of mistake to the Sessions service on its own port today and it
still answers 400. Send it through the router and it is a 200. Same graph,
same mistake, two status codes, decided by which process the request happened
to reach. A client that branches on the status line was right about this graph
until chapter~\ref{ch:enter-the-router} put a router in front of it, and has
been wrong ever since.

\index{breaking change!nothing was removed}%
Nobody removed anything here. The field is still there, still returns the same
connection, and every client that ever called it is refused.

\subsection{The One With No Signal At All}

\index{breaking change!the change that answers correctly}%
The third edit is the one chapter~\ref{ch:null-blast-radius} already made,
applied to one more field: take the exclamation mark off the title of a
session. It composes at exit zero, and the same page query answers:

\begin{minted}[fontsize=\footnotesize,breakanywhere]{json}
{"data":{"sessions":{"nodes":[{"title":"Schemas That Outlive Their Authors"},{"title":"Reading a Query Plan"},{"title":"Paging Without Offsets"},{"title":"The Bank That Split Its Graph"}]}}}
\end{minted}

\index{breaking change!byte for byte the old answer}%
Byte for byte the answer the unchanged graph gives, because every row in this
seed has a title. Introspect the field and the schema has moved: \code{title}
comes back as a plain \code{SCALAR} named \code{String}, where \code{id},
\code{startsAt} and \code{durationMinutes} in the same response still come back
\code{NON\_NULL} with the named type underneath.

\index{breaking change!nothing announces it}%
Nothing announces it. Not the composition, not the router when it loads the
config, not the request, not the router's log. A client that generated its
types from the old schema and has not introspected since keeps a guarantee the
graph stopped making, and nothing at all happens until the day some row has a
null title. Then a null arrives where the generated code has no branch, in a
release nobody connected to a schema change made months earlier by a team the
client's authors have never met.

\index{breaking change!the three sit on one scale}%
Three edits, and they are the same kind of thing at three distances. The
removed field fails immediately and loudly enough that somebody will see it.
The required argument fails immediately and could plausibly be caught by
whoever added it, because the query it breaks is the obvious one. The
nullability change fails at a moment nobody can predict, in a client nobody
told, and the only trace is a null.

\index{HTTP 200!every failure in this book, again}%
And all three come back two hundred. The graph has one way to say no to a
client, and it is the same way it says yes.
\index{breaking change!what a client actually receives|)}%
